\begin{remark}[Both constants in the lemma are attained]
\label{rem:ch49-truncation-sharp}
The upper bound is attained by every constant function \(\psi\equiv c\).
The coefficient in the lower bound is sharp in every dimension.  Indeed,
take \(X=\mathbf P^n\), \(\omega=c_1(\cO(1))\), with homogeneous
coordinates \(z_0,\ldots,z_n\), and set
\[
 \psi:=\max\!\left\{\log|z_1|,
              \max_{j\ne1}\bigl(\log|z_j|-1\bigr)\right\}.
\tag{49.6e}\label{eq:ch49-truncation-sharp-fs}
\]
This is the weighted Fubini--Study function in which \(z_1\) has weight
zero and every other coordinate section has weight \(-1\).  Since the
coordinate sections generate \(\cO(1)\), one has
\(\max_j\log|z_j|=0\) on all of \(X^{\An}\).  Comparing each weight with
\(-1\) and with \(0\) therefore gives, pointwise,
\[
 -1=-1+\max_j\log|z_j|\le\psi
 \le\max_j\log|z_j|=0.
\]
At \(v_{\triv}\) every \(\log|z_j|\) is zero, so \(\psi(v_{\triv})=0\)
and the maximum principle gives \(\sup\psi=0\).

Here is the complete energy computation.  Put
\(P=\operatorname{conv}\{0,e_1^\vee,\ldots,e_n^\vee\}\),
\(P_{\mathrm{BJ}}=-P\), and
\[
 f_P(w)=\max\{0,-w_1,\ldots,-w_n\},\qquad u=e_1.
\]
On the toric skeleton, the convex function associated to
\eqref{eq:ch49-truncation-sharp-fs} is
\[
 f_\psi(w)=f_P(w-u)-1
 =\max\{-1,-w_1,-1-w_2,\ldots,-1-w_n\}.
\]
Since the conjugate of \(f_P\) is zero on \(-P\) and \(+\infty\)
off \(-P\), translation of the argument and subtraction of one give
\[
 f_\psi^\vee(m)=1+\langle u,m\rangle=1+m_1
 \quad(m\in-P),
\]
and \(f_\psi^\vee=+\infty\) off \(-P\).  The normalized toric energy
identity, in the precise convention of this chapter, is
\[
 E_\omega(\eta)=-\frac1{V_P}\int_{-P}f_\eta^\vee(m)\,\rd m,
 \qquad V_P:=\vol(P).
\]
The barycenter of \(-P\) is
\(-\frac1{n+1}(1,\ldots,1)\); hence
\[
 E_\omega(\psi)
 =-\frac1{V_P}\int_{-P}(1+m_1)\,\rd m
 =-\left(1-\frac1{n+1}\right)
 =-\frac n{n+1}.
\]
Thus the lower coefficient is attained in every dimension.

For comparison, the two surface assertions quoted later in the chapter
can also be checked locally.  On \(\mathbf P^2\), let \(p\) be a point
and
\(\eta:=\max\{\log|\mathfrak m_p|,-1\}\).  On the compactified
deformation to the normal cone of \(p\), the function \(\eta\) corresponds
to \(H-E\), where \(H\) is the pullback of the hyperplane class and \(E\)
is the exceptional divisor.  Here \(E\simeq\mathbf P^2\), the restriction
of \(H\) to \(E\) is trivial, and
\(\cO_E(E)=\cO_{\mathbf P^2}(-1)\).  Therefore
\[
 H^3=H^2E=HE^2=0,\qquad E^3=1,\qquad V=1.
\]
Consequently
\[
 E_\omega(\eta)=\frac{(H-E)^3}{3V}=-\frac13.
\]
Since \(-1\le\eta\le0\), this is not the lower endpoint \(-2/3\) of the
estimate in dimension two.

Finally, let \(C\subset\mathbf P^2\) be a smooth conic and put
\(\eta_c:=c\max\{\log|\mathcal I_C|,-1\}\).  On the compactified
deformation to the normal cone of \(C\), it corresponds to \(H-cE\).  If
\(\widetilde P\) is the strict transform of the central \(\mathbf P^2\)
and \(F\) its total transform, then
\(F=\widetilde P+E\) and \(E|_{\widetilde P}=C\).  Thus
\[
 HE\widetilde P=\deg C=2,
 \qquad E^2\widetilde P=C^2=4.
\]
The class \(F\) comes from the base, so its restriction to \(E\) is
trivial; and \(H|_E\) comes from the curve \(C\), so \(H^2E=0\).  Hence
\[
 0=HEF=HE\widetilde P+HE^2,
 \qquad
 0=E^2F=E^2\widetilde P+E^3,
\]
which gives \(HE^2=-2\) and \(E^3=-4\).
The restriction to the strict transform of the central \(\mathbf P^2\)
has degree \(1-2c\) on a line, while the restriction to a fibre of
\(E\to C\) has degree \(c\); thus \(\eta_c\) is \(\omega\)-psh for
\(0\le c\le\tfrac12\).  Moreover
\[
 H^3=H^2E=0,\qquad HE^2=-2,\qquad E^3=-4,
\]
and hence
\[
 E_\omega(\eta_c)=\frac{(H-cE)^3}{3}
 =-2c^2+\frac43c^3.
\]
At \(c=\tfrac12\), one has
\(-\tfrac12\le\eta_{1/2}\le0\),
\(\sup\eta_{1/2}=0\), and
\[
 E_\omega(\eta_{1/2})=-\frac13
 =-\frac23\cdot\frac12.
\]
Thus \(\eta_{1/2}\) is another equality case.  If \(C\) is not a
coordinate line, then \(\ord_C\) is not a toric divisorial valuation for
the standard torus, so this equality case is non-toric in the fixed toric
structure.  These two computations are the ones recorded, in expanded
form, in \Cref{ex:ch46-toric-p2,ex:ch49-conic}.
\end{remark}
